\section{History conventions and branch subdivision}\label{ba:codes}

The general box motions now let us choose a recorder and an observer
separately.  We begin with the two frontier recorders in
Section~\ref{sec:bb-frontiers}, which retain history in the
sense of Bennett~\cite{Bennett1973}.  Later we retain the distinct
neighboring-cell guards and persistent-mark tables needed by other
realizations.  A full-domain inverse, not merely agreement along an
initialized run, is the criterion for sharing a compiler.

For these applications let $A$ be the work alphabet with distinguished
blank, $Q$ the states, $q_0$ the initial state, and $H\subseteq Q$
the halting states.  Write $h_n$ for the work head after $n$ steps and
\[
 \delta(q,a)=(q^+(r),b(r),d(r)),\qquad
 r=(q,a)\in\mathcal R=(Q\setminus H)\times A,
 \quad d(r)\in\{-1,0,1\}.
\]
The initial head is zero and the tape has finitely specified nonblank
cells.  Replace missing instructions by unchanged-letter stay moves to
a halt.  If one halt square is needed, merge halting states and redirect
their incoming instructions; no outgoing rule is changed.  This preserves
initial halting as well.  An unreachable halt state can be added when
needed.  These are the same normalizations as in
Section~\ref{sec:bb-frontiers}.

\subsection{Frontier recorders and their notation}

We use the move-first recorder of Lemma~\ref{lem:bb-movefirst}
with its original notation and complete partial domain.
The delayed-head recorder of Lemma~\ref{lem:bb-direction} records the
work displacement until its final return step.  Its table can also be
initialized with the frontier one cell farther to the right.

\begin{lemma}[Six-rule history recorder]\label{ba:six-compiler}
For $r_0=1$ or $r_0=2$, there is a partial table with the two incoming
properties of Lemma~\ref{lem:bb-onecell}, frontier $r_0+n$ at checkpoint
$n$, and checkpoint control recording both the work state and the
preceding work displacement.  Its next checkpoint takes
$2(r_0+n-h_n)+3$ instructions.  The work head moves only on the last
instruction.
\end{lemma}
\begin{proof}\label{ba:six-proof}
Specialize Lemma~\ref{lem:bb-direction} to a singleton passive track
$O$ and delete that coordinate.  Rename its controls
$(M(q,e),G(\tau),A(q,d),B(q,d))$ as
$(M_{q,e},R_\tau,F_{q,d},L_{q,d})$, and its frontier $F$ as $P$.
The full alphabet is
$A\times(\{\bot,P\}\sqcup\mathcal T)\times\{0,1\}$,
where $\mathcal T=(Q\setminus H)\times\{-1,0,1\}\times A$.
Every row and guard is the same under this bijection, so both incoming
properties hold on the entire partial domain, including unused controls.
Only $M_{q,e}$ with $q\in H$ are halting controls.

Initialize $M_{q_0,0}$ with all marks zero, frontier at $r_0$, and
empty history elsewhere.  The case $r_0=1$ is the initialization already
proved in the shared lemma.  For either allowed $r_0$, checkpoint $n$
has frontier $g=r_0+n>h_n$.  Mark the old head, scan right to $g$,
write the record and install the new frontier at $g+1$, then return to
the mark.  All intervening history is nonfrontier, the new frontier
cell is empty, and there is only one mark.  Hence the same rows apply
when $r_0=2$.  There are $g-h_n-1$ continuing right moves,
$g-h_n$ continuing left moves, and four other instructions, giving
$2(g-h_n)+3$.  The last row clears the mark and makes the work
movement; the invariant is restored with frontier $g+1$.  This proves
finite completion, indefinite legality for a nonhalting run, and the
stated halt equivalence, including initial halting.
\end{proof}

\subsection{From symbolic cylinders to closed rectangles}

Radix encodings turn a push or pop of a tape letter into an affine map;
this is the generalized-shift viewpoint of
Moore~\cite{Moore1990,Moore1991}.  The following calculation adds the
separation of full closed rectangles needed for smooth localization.

\begin{lemma}[Separated rectangle codes]\label{ba:rectangles}
Let a deterministic partial tape table on an alphabet $\Gamma$ of size
$m$ have the two incoming properties of Lemma~\ref{lem:bb-onecell}.
Assign distinct odd digits $e(a)\in\{1,3,\ldots,2m-1\}$ and choose
an integer $B\ge2m+1$.  Its local rules admit effective positive diagonal
affine maps of determinant one between finite families of closed planar
rectangles.  Sources are pairwise separated, as are targets.  Both the
minimal subdivision and the subdivision by both adjacent letters given
below realize every tape in their full symbolic cylinders.
\end{lemma}
\begin{proof}\label{ba:rectangles-proof}
Use Lemma~\ref{lem:bb-onecell} with digits $d(a)=e(a)$.  Its digit
hypotheses hold for every $B\ge2m+1$.  In the notation
\[
 C(a_1a_2\cdots)=\sum_{j\ge1}e(a_j)B^{-j},\quad
 h_a(v)=\frac{e(a)+v}{B},\quad
 I_{a_1\cdots a_l}=h_{a_1}\cdots h_{a_l}([0,1]),
\]
the left stream starts immediately left of the head and the right stream
starts at the head.  Place state $s$ in $k_s+a[0,1]^2$ with a common
rational $a>0$.  The minimal branch maps are exactly
\eqref{eq:bb-one-cell-branches}; scaling both charts by $a$ leaves
$\operatorname{diag}(B^{-d},B^d)$ unchanged.

For the additional full subdivision, restrict the right- and stay-move
sources to $x\in I_\gamma$, for every full letter $\gamma$.
Their targets become $I_{\beta\gamma}\times[0,1]$ and
$I_\gamma\times I_\beta$, respectively, for written letter $\beta$.
The already split left-move targets remain
$[0,1]\times I_{\gamma\beta}$.  At a fixed target state the
incoming direction is fixed, and the pair $(\beta,\gamma)$
distinguishes the branches in every case.  Thus the added subdivision
preserves separate target separation, while source separation follows
by restriction of the shared source rectangles.  All these statements
concern the filled closed rectangles and arbitrary symbolic tails.

For completeness the quantitative gaps remain valid for the entire
range $B\ge2m+1$: distinct first-digit intervals have gap at least
$B^{-1}$ and distinct second-digit intervals at least $B^{-2}$.
The source and target families therefore have within-state gaps at
least $aB^{-2}$; different state squares have their own chosen gaps.
Tail values lie between $1/(B-1)$ and $(2m-1)/(B-1)<1$.
An initial constant tail after $l$ digits contributes
$B^{-l}e(a_{\rm tail})/(B-1)$, so initialized codes are rational.
At checkpoints the history records, or the known absolute origin of
the record block, also recover the absolute work-head position.
\end{proof}


\subsection{Conventions for the remaining geometric realizations}\label{bc:section}
\label{bc:frontier-section}
Unless a different table is specified, subsequent routes use the move-first recorder of
Lemma~\ref{lem:bb-movefirst}, with initial frontier two. The controls
$(C_q,J_r,S_r,P_q,T_q)$ there are written $(W_q,P_r,B_r,C_q,D_q)$
here, and its empty/frontier/record symbols $(e,F,r)$ are written
$(E,F,R_r)$. Permuting work, history and mark coordinates in every row
preserves every guard. This identifies the full partial table, its inverse
and the exact checkpoint count $2(2+n-h')+4$, including arbitrary allowed
tapes. A nonterminal branch means one with both endpoint controls
nonterminal; a final branch into a halt can enter the observer.

The full branches are those of Lemma~\ref{lem:bb-onecell}. For an alphabet
of size $m$ we use one of the following explicitly indicated conventions:
\begin{equation}\label{bc:digits}
(B,\{d_a\})=(2m+2,\{1,3,\ldots,2m-1\}),\quad
(2m+1,\{1,3,\ldots,2m-1\}),\quad
(2m,\{0,2,\ldots,2m-2\}).
\end{equation}
All satisfy its digit bound and separation hypotheses. In the last
convention the full blank has digit zero; boundary codes remain in the
closed-rectangle and neighborhood conclusions. Write
$C(a_1a_2\cdots)=\sum_{j\ge1}d_{a_j}B^{-j}$ and
$I_v=\sum_{j=1}^{|v|}d_{v_j}B^{-j}+[0,B^{-|v|}]$.
A constant tail after $k$ positions contributes
$d_{\rm tail}B^{-k}/(B-1)$, so the initial codes are rational.
Guarded three-cell rules use Lemma~\ref{lem:bb-window-rectangles} with
$(a,b)=(1,2)$ after their own full-domain inverse has been proved.

\subsection{Analytic conventions for the geometric variants}\label{bc:analytic-section}
The whole-space baseline in these variants means smooth
$u\in CH^2\cap C^1L^2$, bounded $u,\nabla u$ on finite intervals,
and a pressure representative in $CH^1$. When the competitor-gradient
condition is omitted explicitly, the other conditions remain unchanged.
Both cases fall under Lemma~\ref{lem:b-comparison}(i); separately stated
pressure classes are retained. The torus baseline is its classical torus
clause. Unless a projection is specified, the force is
\begin{equation}\label{bc:residual}
f_U=U_t+(U\cdot\nabla)U-\nu\Delta U,\qquad (u,p)=(U,0).
\end{equation}
Torus projection is the separate choice in
Proposition~\ref{prop:projection-l2}, with changed pressure.
Localized affine motion is supplied by Lemma~\ref{lem:bb-curl} on a
compact time interval, using effective derivative bounds, positive lower
scale bounds and the positive cutoff margins specified in each route.
The potential is \eqref{eq:bb-potential}; the resulting flow fixes every
region disjoint from its support. In particular exponential diagonal
scales and linear reciprocal scales have respectively generators
$\theta'\log D$ and $\diag(l'/l,-l'/l,0)$.

\begin{lemma}[Finite forcing mechanisms]\label{bc:analytic}
A finite family of the preceding pulses, joined with flat time collars,
can be used once for initialization and repeated with period one thereafter.
It gives a smooth effective velocity and residual force with all mixed
derivatives uniformly bounded. In $\R^3$ both have one fixed compact
spatial support, and the velocity has uniformly bounded kinetic energy.
In a torus chart,
periodized curl or planar Hamiltonian pulses have zero spatial mean, as
does their residual force. Initial velocity and pressure are zero.
Uniqueness holds in the comparison classes of Section~\ref{sec:model};
in the whole-space class the competitor's gradient bound can be omitted
if the remaining conditions are retained.
\end{lemma}
\begin{proof}\label{bc:analytic-proof}
Each pulse is a finite expression in rational data, effective smooth
switches, their derivatives, positive diagonal scales and logarithms.
Finite positive separation margins bound every cutoff derivative. A finite
maximum bounds a period and the loading interval, and repetitions preserve
that bound. To evaluate near a join, a coarse enclosure of time gives a
finite superset of possibly active pulses. Summing them avoids an exact
comparison of a computable time with an integer. The flat collars give
smoothness. Compact support gives every required spatial Sobolev bound.
Lemma~\ref{lem:p2-realization} gives the asserted solution and comparison;
on a torus the integral of each curl or Hamiltonian derivative vanishes,
and the convection term is a divergence. Its flow satisfies
$\det D\Phi_t=1$, because differentiating this determinant gives
$(\divg U)(t,\Phi_t)\det D\Phi_t=0$ and its initial value is one.
The velocity and acceleration identities in that lemma consequently apply
to these volume-preserving diffeomorphisms.
\end{proof}

For later use, spatial chart changes are always made at the velocity level.
If $x=o+a y$, set $\widetilde U(t,x)=aU(t,(x-o)/a)$ and then compute
\eqref{bc:residual} at the prescribed physical $\nu$. The time and convection
terms scale by $a$, while the Laplacian term scales by $a^{-1}$.
Thus chart rescaling does not silently change the fixed physical viscosity.
